\section{Discussion}

The experiments performed in this paper illustrate several properties of the proposed framework. First of all, the propagation algorithm has been shown to stabilize at some credibility distribution and this distribution does not depend on the way sources are initialized. Consequently, one could reasonably assume that credibility estimates depend on the credibility graph structure and not on the initial credibility allocation.

Another property worth noting is the role of graph connectivity. Disconnected sources evolve independently; therefore, their credibility estimates do not propagate through the whole network. It is clear that this consideration called for a new approach to modeling claims. In the experiments, the definition of claims changed to entities and attributes instead of single value which increased graph connectivity and allowed conflicting assertions to be treated within the same inference process.

Furthermore, agreement-weighted propagation showed that the propagation algorithm allows for more information to be used during the propagation process without any additional evaluation. The propagation algorithm can use the level of local agreement between conflicting values as the source of information.

Nonetheless, there are still some problems related to the framework that deserve further research. Graph-based credibility inference shares similarities with other iterative graph ranking methods such as PageRank \cite{brin1998anatomy} and HITS \cite{kleinberg1999authoritative}, as well as truth discovery methods including TruthFinder \cite{yin2008truthfinder} and subsequent surveys of truth discovery \cite{li2016survey}. However, one of the main problems in credibility inference is distinguishing genuine independent evidence from repeated information. The fact that many sources agree upon something does not automatically mean that the consensus is independent because information may propagate due to copying and syndication. Hence, the agreement overestimates the support for a particular claim considerably. Explicitly modeling source dependencies is one of the most important directions for future graph-based credibility inference \cite{dong2009integrating}.